\section*{Limitations}

\paragraph{Models.} We test three generators: two closed OpenAI models and one 7B open-weights model, run at 4-bit
precision with one seed. Larger open models and reasoning models may respond differently, especially to the
directional header. The LLM detector is a single prompt to one model. We compare against no trained detector such as
ConflictRAG's or TrustRAG's filter \citep{wang2026conflictrag,zhou2025trustrag}, which would sit at different points
on the grid.

\paragraph{Data and power.} The confirmatory evidence comes from one benchmark, RAMDocs: 300 questions, 139 of them
conflict questions. Its misinformation is synthetic, and its document labels are imperfect. We exclude 35 hold-out
questions whose status depends on documents we could not label without human annotation, and we check that the
conclusions survive treating those documents as misinformation. Our comparisons detect effects of about
\mdeMed\ points with 80\% power, so smaller effects, including small harms on clean questions, may be missed.
WikiContradict has genuine conflicts but no misinformation side. CONFLICTS, the most realistic benchmark we use, shows
no significant effect for any condition. It was run with one seed, scored only on items with a reference answer, and
without the planned LLM judge, so it can neither confirm nor refute the RAMDocs findings. We do not know how the
findings transfer to naturally occurring misinformation.

\paragraph{Metric.} RAMDocs ships no scorer. Our primary metric reconstructs the published definition and scores our
MADAM-RAG port close to the published number on the pilot split, but the published subset is unknown, so this is not a
reproduction. The metric is string matching: it misses paraphrases, and partial answers interact with ``no wrong
answer''. We report a more lenient variant throughout and mark any conclusion that depends on the choice.

\paragraph{Direction and the grid.} Oracle direction is close to giving the model an answer key, so the oracle rows
are upper bounds. In the synthetic grid only perfect direction gives a significant overall gain. The real LLM
detector, at \dirAccLlm\% direction accuracy, gains on conflict questions over a neutral flag, but its net gain over no flags
is significant only with gpt-4.1. Synthetic false flags are spread uniformly and their suspect side is random, whereas
real detectors concentrate their errors. The grid predicts real detectors only roughly, and three detectors cannot
show whether document- or pair-level summaries predict better.

\paragraph{Scope.} We study flags passed to the generator as text. Learned interfaces such as soft prompts or
credibility fine-tuning \citep{ye2026tcr,pan2024credibility}, and retrieval-side filtering, may follow different
trade-offs. We test one adversarial case, a detector that always names the wrong side, but not an attacker who controls
which pairs are flagged.

\section*{Ethics Statement}
This work uses public benchmarks and commercial model APIs, and involves no human subjects or personal data. Its main
risk is dual use. Our results show that a false ``likely misinformation'' label, or the removal of a passage, can
change a RAG system's answers: an inverted oracle that names the correct side as misinformation costs up to
\aadvRemConf\ points. An adversary who can trigger a deployed detector could exploit this
\citep{zou2025poisonedrag}. We therefore report harm on conflict-free questions alongside every gain, and we argue that
systems should measure what a detector does at the document level before acting on its flags.
